\documentclass[10pt,a4paper]{amsart}
\usepackage[margin=19mm]{geometry}
\usepackage{amsmath,amssymb,mathtools}
\usepackage[scr=rsfs,cal=euler]{mathalfa}
\usepackage{xfrac}
\usepackage[unicode,hidelinks,bookmarks=false]{hyperref}
\newcommand{\ie}{i.e.\ }
\newcommand{\cf}{cf.\ }
\newcommand{\resp}{respectively\ }
\newcommand{\Subsection}{\subsection}
\providecommand{\nobreakdash}{\nobreak\hbox{-}}
\setlength{\emergencystretch}{2em}
\multlinegap0pt
\usepackage{amssymb}
\newtheorem{theorem}{Theorem}
\newcommand{\de}{\delta}
\newcommand{\si}{\sigma}
\title{Evaluation of iterated Ore polynomials and skew Reed-Muller codes}
\author{Andre Leroy, Nabil Bennenni}
\date{Source: arXiv:2603.01287v1 (CC BY 4.0)}
\begin{document}
\maketitle

\noindent\emph{Source and licence.} Evaluation of iterated Ore polynomials and skew Reed-Muller codes. Version arXiv:2603.01287v1 (CC BY 4.0), distributed by arXiv under the Creative Commons Attribution 4.0 International licence, http://creativecommons.org/licenses/by/4.0/, as recorded in the arXiv licence metadata for this exact version and shown on the abstract page https://arxiv.org/abs/2603.01287v1. Full licence notice: \texttt{licenses/arxiv-2603.01287-CC-BY-4.0.txt}.\par\smallskip

\noindent\textit{Editorial scope.} Counted unit: the article's theorem fields (source0.tex lines 451--468). The article's complete original proof of it is delivered with it (source0.tex lines 469--511, one \texttt{proof} environment of 43 lines). No mathematics is written, repaired or re-derived: the statement and the proof are the article's own lines, in the article's own notation, and no retained line is substituted or re-flowed.  Retained uncounted as the source-local context the proof needs: lines 428--442.  Nothing was omitted from any retained span, so the package carries no omission record.  No retained line contains a citation, so the wrapper carries no bibliography and no bibliography entry.  No retained line contains a figure environment, an image-inclusion command or drawing code, so no image is delivered and the package declares no asset dependency.  Editorial formatting only, in addition: an AMS \texttt{amsart} preamble that restates the article's own declarations and macros used by the retained lines, namely the environments \texttt{theorem} on separate counters, together with the standard packages those lines need.  The retained environments print in this document as Theorem 1 in order of appearance, whereas the article prints the same items with the numbering of its own full document.  The complete native source of the article as distributed in the arXiv e-print for this exact version is bundled unchanged as \texttt{sources/195550/source0.tex}.
\begin{theorem}
	\label{generalization of Gordon-Motzkin}
	
	Let $a$ be an element of a division ring $K$ and $f(t)\in R=K[t;\sigma,\delta]$ be a polynomial of degree $n$.
	Then:
	\begin{enumerate}
		\item[(1)] $C^{\si,\de}(a)$ is a subdivision ring of $K$.
		\item[(2)] The map $T_a$ is a right $C^{\si,\de}(a)$ linear map.
		\item[(3)] $f(t)$ has roots in at most $n$ $(\sigma,
		\delta)$-conjugacy classes, say $\{\Delta
		(a_1),\dots,\Delta(a_r)\}$, $r\le n$;
		\item[(4)] $\sum_{i=1}^rdim_{C(a_i)}\ker (f(T_{a_i}))\le n$, where 
		$C(a_i):=C^{\si,\de}(a_i)$ for $1\le i \le r$.
	\end{enumerate}
\end{theorem}

\begin{theorem}
	\label{finite fields}	
	Let $p$ be a prime number and $\mathbb F_q$ be the finite field
	with $q=p^n$ elements.  Denote by $\theta$ the Frobenius
	automorphism.  Then:
	\begin{enumerate}
		\item[a)] There are $p$ distinct $\theta$-conjugacy classes in $\mathbb F_q$.
		\item[b)] $C^{\theta}(0)=\mathbb F_q$ and, for $0\ne a\in \mathbb F_q$, we have
		$C^{\theta}(a)=\mathbb F_p$.
		\item[c)] In $\mathbb F_q[t;\theta]$, the least left common multiple of all the elements
		of the form $t-a$ for $ a\in\mathbb F_q$ is the polynomial
		$G(t):=t^{(p-1)n +1}-t$.  In other words, $G(t)\in \mathbb
		F_q[t;\theta]$ is of minimal degree such that $G(a)=0$ for all
		$a\in \mathbb F_q$.
		\item[d)] The polynomial $G(t)$ obtained in c) above is invariant,
		i.e. $RG(t)=G(t)R$.
	\end{enumerate}
\end{theorem}

\begin{proof}
	a)   Let us denote by $g$ a generator of the cyclic group $\mathbb
	F_q^{*}:=\mathbb F_q \setminus\{0\}$. The $\theta$-conjugacy class
	determined by the zero element is reduced to $\{0\}$ i.e. $\Delta (0)=\{0\}$.
	The $\theta$-conjugacy class determined by $1$ is a subgroup of $\mathbb
	F_q^*$: $\Delta (1)=\{\theta (x)x^{-1}\,|\, 0\ne x\in \mathbb F_q
	\}=\{x^{p-1}\,|\, 0\ne x\in \mathbb F_q\}$.  It is easy to check
	that $\Delta (1)$ is cyclic generated by $g^{p-1}$ and has order
	$\frac{p^n-1}{p-1}$.  Its index is $(\mathbb F_q^* : \Delta
	(1))=p-1$.   Since two nonzero elements $a,b$ are
	$\theta$-conjugate if and only if $ab^{-1}\in \Delta (1)$,  we
	indeed get that the number of different nonzero $\theta$-conjugacy classes
	is $p-1$.   This yields the result.
	
	\noindent b) If $a\in \mathbb F_q$ is nonzero, then
	$C^{\theta}(a)=\{x\in \mathbb F_q\,|\, \theta(x)a=ax\}$ i.e.
	$C^{\theta}(a)=\mathbb F_p$.
	
	\noindent c)  We have, for any $x\in \mathbb F_q,\,
	(t^{(p-1)n+1}-t)(x)=\theta^{(p-1)n}(x)\dots \theta(x)x - x$. Since
	$\theta^n=id$, and $N_n(x):=\theta^{n-1}(x)\dots
	\theta(x)x \in \mathbb
	F_p$, we get $(t^{(p-1)n+1}-t)(x)=
	x(\theta^{n-1}(x)\dots\theta(x)x )^{p-1}-x=
	xN_n(x)^{p-1}-x=0$.  This shows that indeed $G(t)$ 
	annihilates all the elements
	of $\mathbb F_q$ and hence $G(t)$ is a  left common multiple of
	the linear polynomials $\{(t-a)\,|\,a\in \mathbb F_q\}$.  Let $h(t):=[t-a\,|\,a\in \mathbb F_q]_l$ denote their least left common multiple. It remains to show that
	$\deg h(t)\ge n(p-1)+1$.  Let $0=a_0,a_1,\dots,a_{p-1}\in \mathbb{F}_q$ be
	elements  representing the $\theta$-conjugacy classes (Cf. a) above).  Denote by $C_0,C_1,\dots,C_{p-1}$ their respective
	$\theta$-centralizer.  The formulas recalled in the paragraph before  Theorem \ref{generalization of Gordon-Motzkin} shows that 
	$h(T_a)(x)=h(a^x)x=0$ for any nonzero element $x\in \mathbb F_q$ and any element 
	$a\in\{a_0,\dots,a_{p-1}\}$.   
	Hence $\ker h(T_{a_i})=\mathbb F_q$ for $0\le i \le p-1$. Using
	Inequality $(4)$ in Theorem \ref{generalization of
		Gordon-Motzkin} and the statement $b)$ above, we get $\deg h(t)\ge
	\sum_{i=0}^{p-1} dim _{C_i}\ker h(T_{a_i})=dim_{\mathbb F_q}\mathbb
	F_q + \sum_{i=1}^{p-1}dim_{\mathbb F_p}\mathbb F_q =1+(p-1)n$, as
	required.
	
	\noindent d)  Since $\theta^n=id.$, we have immediately that
	$G(t)x=\theta(x)G(t)$ and obviously $G(t)t=tG(t)$.
\end{proof}

\end{document}
